\begin{man}{rand,srand}
    \begin{synopsis}
        \begin{plaincode}[gobble=8]
            #include <stdlib.h>

            int rand();
            void srand(unsigned int seed);
        \end{plaincode}
    \end{synopsis}

    \begin{desc}
        Der Befehl \texttt{rand()} generiert eine Zufallszahl von 0 bis zum dem im Makro \texttt{RAND\_MAX} definierten Wert. Diese Zufahlszahl wird mit einem Algorithmus berechnet und hängt von ihrem Anfangswert ab (das ist die ``seed'', auf Deutsch ``Saat''). Dieser Anfangswert kann mit dem Befehl \texttt{srand} geändert werden.
    \end{desc}

    \begin{example}
        \begin{code}
#include <stdio.h>
#include <stdlib.h>
#include <time.h>

int main() {
    srand(time(NULL));
    printf("Zufallszahl von %d bis %d: %d\n", 0, RAND_MAX, rand());
    return 0;
}
        \end{code}
    \end{example}

    % \begin{returns}
    % \end{returns}

    \begin{seeother}
        \manref{time}
    \end{seeother}
\end{man}

